\usepackage{blindtext}

% Axiomathic web export. lwarp must load before array/tikz/etc., but
% standalone subfiles should remain ordinary PDFs.
\makeatletter
\@ifclassloaded{subfiles}{}{%%
  \usepackage[mathjax,HomeHTMLFilename=index]{lwarp}
  \setcounter{FileDepth}{0}
  \setcounter{SideTOCDepth}{1}
  \CSSFilename{site.css}
  \providecommand{\cs}[1]{\texttt{\textbackslash#1}}
}
\makeatother


% Graphics / maths
\usepackage{graphicx}
\usepackage{amsmath}
\usepackage{amssymb}
\usepackage{amsthm}

% Theorem styling
\usepackage{xcolor}
\usepackage{thmtools}
\usepackage{tcolorbox}
\usepackage{mdframed}
\tcbuselibrary{theorems}


% Other packages
\usepackage{xparse}
\usepackage{float}
\usepackage{imakeidx}

\usepackage{tabularx}
\usepackage{colortbl}
\usepackage{longtable}
\usepackage{arydshln}
\usepackage{array}
\usepackage{subfiles}

\usepackage{tikz}

%specifically for checklist tables
\newcolumntype{R}{>{\raggedright\arraybackslash}p{3.0cm}}

%completed tickbox symbol
\newcommand{\squaretick}{%
  \makebox[0pt][l]{$\square$}%
  \raisebox{.15ex}{\hspace{0.1em}$\checkmark$}%
}

%tally counter
\newcommand{\Tally}[2][]{\begin{tikzpicture}[baseline,#1]
\foreach \X [evaluate=\X as \Y using {int(mod(\X,5))}]in {1,...,#2}
{\ifnum\Y=0
\draw (\X*0.5ex+0.3ex,0) -- ++(-2.2ex,1.8ex);
\else
\draw (\X*0.5ex+0.3ex,0) -- ++(0.3ex,1.8ex);
\fi}
\end{tikzpicture}}

\def\multiset#1#2{\ensuremath{\left(\kern-.3em\left(\genfrac{}{}{0pt}{}{#1}{#2}\right)\kern-.3em\right)}}
\newmdtheoremenv{boxtheorem}{Theorem}
\newcommand{\textoverline}[1]{$\overline{\mbox{#1}}$}
\makeindex
\DeclareUnicodeCharacter{2009}{ }
\newcommand\mybox[2][]{\tikz[overlay]\node[fill=blue!20,inner sep=2pt, anchor=text, rectangle, rounded corners=0mm,#1] {#2};\phantom{#2}}

% Set the overall layout of the tree
\tikzstyle{level 1}=[level distance=3.5cm, sibling distance=3.5cm]
\tikzstyle{level 2}=[level distance=3.5cm, sibling distance=2cm]
% Define styles for bags and leafs
\tikzstyle{bag} = [text width=4em, text centered]
\tikzstyle{end} = [circle, minimum width=3pt,fill, inner sep=0pt]




%custom theorem-like environments
\declaretheorem[
numberwithin=chapter,  shaded={ padding=6pt, leftmargin=-6pt, rightmargin=-6pt, rulecolor=red!10, rulewidth=2pt, bgcolor=red!10},
  name=Theorem,
]{theorem}

\declaretheorem[
numberlike=theorem,  shaded={ padding=6pt, leftmargin=-6pt, rightmargin=-6pt, rulecolor=yellow!26, rulewidth=0pt, bgcolor=yellow!26},
  name=Lemma,
]{lemma}

\declaretheorem[
numberlike=theorem,  shaded={ padding=6pt, leftmargin=-6pt, rightmargin=-6pt, rulecolor=yellow!26, rulewidth=0pt, bgcolor=yellow!26},
  name=Proposition,
]{prop}

\declaretheorem[
numberlike=theorem,  shaded={ padding=6pt, leftmargin=-6pt, rightmargin=-6pt, rulecolor=yellow!26, rulewidth=0pt, bgcolor=yellow!26},
  name=Corollary,
]{corollary}

\declaretheorem[
numberlike=theorem,  shaded={ padding=6pt, leftmargin=-6pt, rightmargin=-6pt, rulecolor=cyan!18, rulewidth=0pt, bgcolor=cyan!18},
  name=Definition,
]{defn}

\declaretheorem[
numberlike=theorem,  shaded={ padding=6pt, leftmargin=-6pt, rightmargin=-6pt, rulecolor=cyan!18, rulewidth=0pt, bgcolor=cyan!18},
  name=Notation,
]{nota}

\declaretheorem[
numberlike=theorem,  shaded={ padding=6pt, leftmargin=-6pt, rightmargin=-6pt, rulecolor=gray!12, rulewidth=0pt, bgcolor=gray!12},
  name=Conjecture,
]{conj}


\declaretheorem[
numbered=no,  shaded={ padding=6pt, leftmargin=-6pt, rightmargin=-6pt, rulecolor=black, rulewidth=1pt, bgcolor=green!12},
  name=Example,
]{example}

\declaretheorem[
numbered=no,  shaded={ padding=6pt, leftmargin=-6pt, rightmargin=-6pt, rulecolor=black, rulewidth=0pt, bgcolor=white},
  name=Remark,
]{remark}
% Commands required by Shuffling Cards when included as a subfile
\ExplSyntaxOn
\NewDocumentCommand{\cycle}{ O{\;} m }
 {
  (
  \axiom_cycle:nn { #1 } { #2 }
  )
 }
\seq_new:N \l_axiom_cycle_seq
\cs_new_protected:Npn \axiom_cycle:nn #1 #2
 {
  \seq_set_split:Nnn \l_axiom_cycle_seq { , } { #2 }
  \seq_use:Nn \l_axiom_cycle_seq { #1 }
 }
\ExplSyntaxOff
\DeclareMathOperator{\lcm}{lcm}

% Dependencies required by Easier Waring when included as a subfile
\usepackage{mathtools}
\usepackage{booktabs}
\usepackage{enumitem}
\usepackage{microtype}
\usepackage{hyperref}
\usepackage{cleveref}
\usepackage{bm}
\hypersetup{colorlinks=true,linkcolor=blue!55!black,citecolor=blue!55!black,urlcolor=blue!55!black}
\newtheorem{warning}[theorem]{Caution}
\newcommand{\Z}{\mathbb Z}
\newcommand{\N}{\mathbb N}
\newcommand{\Q}{\mathbb Q}
\newcommand{\T}{\mathbb T}
\newcommand{\F}{\mathbb F}
\newcommand{\eps}{\varepsilon}
\newcommand{\e}{\mathrm e}
\newcommand{\calP}{\mathcal P}
\newcommand{\calR}{\mathcal R}
\newcommand{\ord}{\operatorname{ord}}
